\documentclass[11pt,a4paper]{article}
\usepackage{agpstyle}
\usepackage{pgfplots}
\pgfplotsset{compat=1.17}
\usepgfplotslibrary{fillbetween}
\usetikzlibrary{patterns,shapes.geometric}
\usepackage[section]{placeins}   % keep each section's floats inside that section

% Let large figures sit next to their text instead of drifting to the end.
\renewcommand{\topfraction}{0.9}
\renewcommand{\bottomfraction}{0.8}
\renewcommand{\textfraction}{0.07}
\renewcommand{\floatpagefraction}{0.75}
\setcounter{topnumber}{3}
\setcounter{totalnumber}{4}

\title{\bfseries Reinforcement Learning for the Continuous Art Gallery Problem\\[4pt]
\large Preliminaries, Policy Architecture, and MDP Formulation}
\author{AGP Project --- Technical Report}
\date{\today}

% A visible marker for numbers that do not exist yet.
\newcommand{\pending}[1]{\textcolor{agpaccent}{\textbf{[pending#1]}}}

% One colour scheme for every figure.
\colorlet{seen}{agpblue!22}
\colorlet{unseen}{agpaccent!35}
\colorlet{wall}{black!80}
\colorlet{masked}{black!18}
\tikzset{
  hole/.style={fill=white, draw=agpaccent, line width=0.9pt},
  outer/.style={draw=wall, line width=1.1pt},
  guard/.style={circle, fill=black, inner sep=1.6pt},
  ray/.style={dashed, black!55, line width=0.5pt},
  blk/.style={draw=agpblue, fill=agpgrey, rounded corners=2pt, align=center, font=\small, inner sep=4pt},
  flow/.style={-Stealth, thick, agpblue},
}

\begin{document}
\maketitle

\begin{abstract}
\noindent
We train a PPO agent to place point guards in polygonal galleries with holes so that at
least $98\%$ of the free space is seen, using as few guards as possible. This report
explains every design decision in the current implementation. \Cref{sec:prelim} gives the
geometry that motivates a two-stage curriculum: coverage is monotone submodular, so a
guard's redundancy can only be measured \emph{after} the gallery is saturated, and the
dense rewards we use are potential-based, so they telescope into the true objective.
\Cref{sec:arch} describes the policy network and why each part of it exists: it keeps the
map's spatial layout to the output and factors each placement into a choice of region and a
\emph{continuous} point inside it. \Cref{sec:mdp} states the MDP and the reward with the exact
constants used in code. \Cref{sec:stage3} outlines an optional third stage with a relocation
action.
\end{abstract}

\tableofcontents

% ===========================================================================
\section{Introduction}
% ===========================================================================

The agent sees an image of a gallery and the guards placed so far. At each step it either
adds a guard at a point, deletes the guard nearest to a point, or stops. Training runs in
two stages on one network (\cref{fig:curriculum}):

\begin{description}[leftmargin=1.6em]
  \item[\stageI{} (stage 1).] Only $\ADD$ is allowed and guards are free. The agent learns
        to cover arbitrary galleries.
  \item[\stageII{} (stage 2).] $\DEL$ and $\STOP$ are unmasked, every guard costs reward,
        and deleting a guard that costs no coverage is paid. The agent learns to remove the
        guards it does not need.
\end{description}

Three facts shape everything below. The problem is NP-hard, so a learned heuristic is a
reasonable target (\cref{sec:complexity}). Coverage is submodular, which fixes the
\emph{order} of the two stages (\cref{sec:submod}). And guards see without a range limit,
so whether a point is a good guard position depends on the whole gallery, which shapes the
network (\cref{sec:arch}). The implementation lives in \code{agp/}: geometry in
\code{gallery.py}, the MDP in \code{env.py}, rewards in \code{rewards.py}, the policy in
\code{policy.py}, and the curriculum in \code{train.py}. All constants quoted below are the
values in that code.

% ===========================================================================
\section{Preliminaries: why two stages}\label{sec:prelim}
% ===========================================================================

\subsection{Galleries, visibility, coverage}

\begin{definition}[Gallery and free space]
A gallery $\Gallery = (\Poly, \Holes)$ is a simple polygon $\Poly \subset \Real^2$ with a
finite set $\Holes$ of pairwise disjoint polygonal holes inside it. The free space is
$\Free = \Poly \setminus \bigcup_{H \in \Holes} \operatorname{int} H$, with area
$A = \Area{\Free} > 0$.
\end{definition}

\begin{definition}[Visibility polygon]
A guard at $g \in \Free$ sees $\Vis(g) = \{ q \in \Free : \overline{gq} \subseteq \Free \}$.
\end{definition}

Visibility has \emph{unlimited range}: a guard sees every point it has an unobstructed
segment to, however far, and only the outer wall and the holes block it
(\cref{fig:visibility}). \code{Gallery.visibility\_polygon} computes $\Vis(g)$ exactly. It
casts one ray towards every vertex of $\Poly$ and of every hole, plus one ray
$\pm 10^{-6}$\,rad to either side, so that both edges meeting at an occluding (reflex) corner
are captured. Each ray is $3\times$ the bounding-box diagonal long, so effectively infinite,
and is cut at its first boundary hit. The hits, taken in angular order, form $\Vis(g)$,
which is clipped back to $\Free$.

\begin{figure}[!htbp]
\centering
\begin{tikzpicture}[scale=0.95]
  % free space, all visible to start with
  \fill[seen] (0,0) -- (8,0) -- (8,5) -- (2,5) -- (2,3.6) -- (0,3.6) -- cycle;
  % shadow cast by the hole (rays through its two silhouette corners)
  \fill[unseen] (4,3) -- (8,3.8) -- (8,1.2) -- (4,2) -- cycle;
  % wedge hidden behind the reflex corner of the outer wall
  \fill[unseen] (2,3.6) -- (2,5) -- (2.636,5) -- cycle;
  \draw[ray] (1.5,2.5) -- (8,3.8);
  \draw[ray] (1.5,2.5) -- (8,1.2);
  \draw[ray] (1.5,2.5) -- (2.636,5);
  \draw[outer] (0,0) -- (8,0) -- (8,5) -- (2,5) -- (2,3.6) -- (0,3.6) -- cycle;
  \draw[hole] (4,2) rectangle (5,3);
  \node[guard, label={[font=\small]below:$g$}] at (1.5,2.5) {};
  \node[font=\small] at (4.5,2.5) {hole};
  \node[font=\small, align=center] at (6.7,2.5) {shadow\\of the hole};
  \node[font=\small, anchor=west] (rc) at (3.2,5.55) {hidden behind a reflex corner};
  \draw[-Stealth, black!60] (rc.west) -- (2.3,4.85);
  \node[font=\small, anchor=west] (vs) at (8.4,0.6) {$\Vis(g)$};
  \draw[-Stealth, black!60] (vs.west) -- (6.5,0.5);
\end{tikzpicture}
\caption{\textbf{Visibility has unlimited range and is blocked only by walls and holes.} The
guard $g$ sees the whole blue region, up to the far wall. It cannot see the orange regions:
the shadow cast by the hole, bounded by the rays through the hole's two silhouette corners,
and a thin wedge behind the reflex corner of the outer wall. Thin wedges like this are what
the last few guards of a solution exist to cover.}
\label{fig:visibility}
\end{figure}

\begin{definition}[Coverage]
For a finite guard set $S \subset \Free$ let $\Cov(S) = \bigcup_{g\in S}\Vis(g)$ and
$f(S) = \Area{\Cov(S)}/A \in [0,1]$. A gallery is \emph{solved} when $f(S) \ge c^\dagger$
with $c^\dagger = 0.98$.
\end{definition}

The threshold $c^\dagger < 1$ is deliberate. Full coverage can require a guard for a
sliver of a few square units behind a hole, and that last guard dominates the count
without telling us anything about placement quality. At $98\%$ the greedy baseline solves
$598/600$ training and $88/90$ test galleries under the 16-guard budget; the rest stop
between $97.7\%$ and $98\%$ because no candidate point adds new area
(\cref{sec:greedy}).

\paragraph{Two representations of the same geometry.}
Every quantity that enters the \emph{reward} is computed exactly with Shapely polygons:
$f(S)$, the area a new guard adds, and the area lost when a guard is removed. The
\emph{observation} is a $64\times 64$ raster of the gallery's bounding box: pixel $(r,c)$
samples the point $(x_{\min} + c\,w/63,\; y_{\min} + r\,h/63)$, where $w, h$ are the box's
width and height. The raster is also used for one cheap estimate shown to the agent, the
per-guard exclusive area (\cref{sec:obs}). Exact rewards mean the agent is never paid for
raster artefacts; a raster observation means the network can be convolutional.

\subsection{Complexity, and what a learned policy can gain}\label{sec:complexity}

Chv\'atal's theorem gives $\lfloor n/3\rfloor$ vertex guards as the tight worst case for a
simple $n$-gon \cite{chvatal1975,fisk1978}. Finding a \emph{minimum} guard set is NP-hard
\cite{leelin1986}, hard to approximate \cite{eidenbenz2001}, and with guards allowed
anywhere in the polygon the decision problem is $\exists\Real$-complete
\cite{abrahamsen2018}. Exact and near-exact solvers exist but restrict guards to finite
candidate sets such as vertices or grid points \cite{couto2011,tozoni2016}.

This tells us where a learned policy can do better than classical heuristics:
\begin{enumerate}
  \item \textbf{Continuous positions.} Optimal point guards need not lie on any candidate
        set. A policy that outputs real coordinates can reach positions no candidate-based
        method can. We keep the action continuous for this reason (\cref{sec:action}).
  \item \textbf{Pruning.} Greedy covers are not minimal, so removing redundant guards
        recovers slack that greedy leaves behind (\cref{sec:greedy}).
\end{enumerate}

\subsection{Submodularity: why pruning must come second}\label{sec:submod}

\begin{proposition}[Monotone submodularity]\label{prop:submod}
$f(\emptyset)=0$, $f$ is monotone, and for all $S \subseteq S'$ and every $g$,
\[
  f(S\cup\{g\}) - f(S) \;\ge\; f(S'\cup\{g\}) - f(S').
\]
\end{proposition}

\begin{proof}
With $\nu = \Area{\cdot}/A$, the marginal gain is $\nu(\Vis(g)\setminus\Cov(S))$. Since
$\Cov(S)\subseteq\Cov(S')$, the set being measured can only shrink as $S$ grows.
\end{proof}

Two quantities drive the design. For a candidate $g \notin S$ and a placed guard $g_j\in S$,
\[
  \delta^{+}(g\mid S) = f(S\cup\{g\}) - f(S),
  \qquad
  \delta^{-}(g_j\mid S) = f(S) - f(S\setminus\{g_j\}).
\]
$\delta^+$ is the \emph{marginal gain} of adding, and $\delta^-$ is the \emph{exclusive area}
of a placed guard, i.e.\ what deleting it would cost. A guard with $\delta^-=0$ is strictly
redundant.

\begin{proposition}[Redundancy is revealed only after saturation]\label{prop:redundancy}
Insert guards $g_1,\dots,g_K$ in order and let $S_j=\{g_1,\dots,g_j\}$. Then for every $j$,
\[
  \delta^-(g_j\mid S_K) \;\le\; \delta^+(g_j\mid S_{j-1}),
\]
and $\delta^-(g_j\mid S)$ is non-increasing in $S \ni g_j$.
\end{proposition}

\begin{proof}
Apply \cref{prop:submod} to $S_{j-1}\subseteq S_K\setminus\{g_j\}$.
\end{proof}

\Cref{fig:submod} shows both propositions on a real test gallery. The coverage curve flattens
as guards are added, and the second guard, which added $28\%$ of the gallery when it was
placed, ends up with an exclusive area of zero.

\begin{figure}[!htbp]
\centering
\begin{tikzpicture}
\begin{axis}[
  name=cov, width=0.47\linewidth, height=5.2cm,
  xlabel={guards placed by greedy}, ylabel={coverage $f(S_j)$},
  xmin=0, xmax=7.3, ymin=0, ymax=1.05, xtick={0,...,7},
  ytick={0,0.2,0.4,0.6,0.8,0.98}, yticklabels={0,0.2,0.4,0.6,0.8,$c^\dagger$},
  grid=major, grid style={black!10}, tick label style={font=\small}, label style={font=\small}]
  \addplot[agpblue, thick, mark=*, mark size=1.6pt] coordinates
    {(0,0) (1,0.4259) (2,0.7073) (3,0.8024) (4,0.8858) (5,0.9468) (6,0.9783) (7,0.9859)};
  \addplot[agpaccent, dashed, thick, domain=0:7.3] {0.98};
\end{axis}
\begin{axis}[
  at={(cov.east)}, anchor=west, xshift=1.6cm,
  width=0.47\linewidth, height=5.2cm, ybar, bar width=5pt,
  xlabel={guard $j$ (insertion order)}, ylabel={fraction of gallery},
  xmin=0.4, xmax=7.6, ymin=0, ymax=0.47, xtick={1,...,7},
  legend style={font=\small, draw=none, fill=none, at={(0.98,0.97)}, anchor=north east},
  tick label style={font=\small}, label style={font=\small}]
  \addplot[fill=agpblue!70, draw=agpblue] coordinates
    {(1,0.4259) (2,0.2814) (3,0.0950) (4,0.0835) (5,0.0610) (6,0.0315) (7,0.0076)};
  \addplot[fill=agpaccent!70, draw=agpaccent] coordinates
    {(1,0.1453) (2,0.0000) (3,0.0858) (4,0.0546) (5,0.0491) (6,0.0119) (7,0.0076)};
  \legend{$\delta^+$ when placed, $\delta^-$ at the end}
  \node[font=\small, align=center, anchor=south] at (axis cs:2.35,0.29) {redundant\\at the end};
\end{axis}
\end{tikzpicture}
\caption{\textbf{Greedy saturates with diminishing returns and leaves redundant guards behind.}
Real data: greedy cover of hard test gallery \#63. \emph{Left:} each extra guard adds less
(\cref{prop:submod}); the seventh adds $0.8\%$ to cross $c^\dagger$. \emph{Right:} for every
guard, the exclusive area at the end ($\delta^-$) is at most its gain when placed
($\delta^+$), as \cref{prop:redundancy} requires. Guard 2 covered $28\%$ when placed and is
fully redundant at the end, so the greedy prune removes it ($k_{\mathrm{greedy}}=7$,
$k_{\mathrm{ref}}=6$). Its redundancy only became visible after guards 3--7 were placed.}
\label{fig:submod}
\end{figure}

\takeaway{%
\textbf{Consequence.} A guard can only become \emph{cheaper} to delete as more guards are
added. A deletion decision made early in an episode sees an exclusive area that is too large;
after saturation it sees the smallest value that guard will ever have. Learning to cover
first and to prune second is therefore the order in which the pruning signal is most
informative, not just a convenience.}

\subsection{Greedy, its guarantees, and the reference count}\label{sec:greedy}

The greedy rule ``add the guard of largest $\delta^+$'' inherits the $(1-1/e)$ guarantee for
budgeted monotone submodular maximisation \cite{nemhauser1978} and the logarithmic-factor
guarantee for full cover, which is essentially the best possible for set cover in polynomial
time \cite{feige1998}. Greedy is therefore strong at \emph{covering}. It is weak at
\emph{minimality}: by \cref{prop:redundancy}, guards chosen early lose exclusive area as
later guards overlap them, and greedy never revisits a choice.

Our baseline (\code{baselines.py}) makes this concrete for each gallery:
\begin{enumerate}
  \item \textbf{Candidates.} One point per vertex of $\Poly$ and of every hole, moved $1\%$
        of the diagonal towards an interior reference point, plus $40$ uniform interior
        samples.
  \item \textbf{Greedy cover.} Repeatedly add the candidate with the largest raster marginal
        gain until $f\ge c^\dagger$, the budget of $16$ is reached, or no candidate adds
        area. The count is $k_{\mathrm{greedy}}$.
  \item \textbf{Greedy prune.} Repeatedly delete the guard with the smallest exact
        $\delta^-$ while coverage stays $\ge c^\dagger$. The remaining count is
        $k_{\mathrm{ref}}$.
\end{enumerate}

$k_{\mathrm{ref}}$ is the per-instance reference that the parsimony reward and the reported
\emph{guard ratio} $k/k_{\mathrm{ref}}$ are measured against. It is an \emph{upper bound} on
the optimum, not the optimum: it depends on the candidate set, and the true minimum is
NP-hard to compute. A guard ratio below $1$ means beating greedy, not proving optimality.
On the training set, $k_{\mathrm{greedy}}$ ranges over $1$--$3$ (easy), $2$--$5$ (mid) and
$3$--$11$ (hard). Mean test-set values are in \cref{tab:greedy}.

\begin{table}[!htbp]
\centering\small
\caption{Greedy baseline on the 90 test galleries (30 per tier).}
\label{tab:greedy}
\begin{tabular}{lccc}
\toprule
tier & mean $k_{\mathrm{greedy}}$ & mean $k_{\mathrm{ref}}$ (after prune) & mean coverage \\
\midrule
easy & 1.93 & 1.93 & 0.994 \\
mid  & 2.70 & 2.67 & 0.992 \\
hard & 6.07 & 5.97 & 0.986 \\
\bottomrule
\end{tabular}
\end{table}

\subsection{Why one stage collapses}\label{sec:collapse}

The objective is to solve the gallery with few guards. With a guard cost $w_{\mathrm{par}}$
per guard and a coverage weight $w_{\mathrm{cov}}$, the immediate value of one more $\ADD$ is
$w_{\mathrm{cov}}\,\delta^+ - w_{\mathrm{par}}$, which is negative whenever
\[
  \delta^+ \;<\; w_{\mathrm{par}}/w_{\mathrm{cov}} \;=\; 0.15 \quad\text{(our stage-2 weights).}
\]
In \cref{fig:submod}, guards 3--7 all add less than $15\%$. Their only justification is the
terminal bonus for crossing $c^\dagger$, and a policy that cannot yet cross $c^\dagger$
reliably never collects it. From such a policy's point of view the late guards are pure
cost, so it learns to place fewer of them, which makes crossing $c^\dagger$ rarer still. The
fixed point of this feedback is ``place few guards and fail''.

Stage 1 breaks the loop by setting $w_{\mathrm{par}}=0$: every $\ADD$ has non-negative value,
the goal bonus is reachable, and the agent learns to cover. Stage 2 turns the cost on only
after covering is reliable (\cref{fig:curriculum}). This matches the reference locomotion
work, where a single-stage dribbling policy reaches $33\%$ success against $86.7\%$ with two
stages \cite{wang2025dribble}.

\subsection{Potential-based shaping: dense reward without bias}\label{sec:shaping}

PPO needs a dense signal, but dense rewards can change which policy is optimal, and a
reward paid per \emph{event} can be farmed by repeating the event. We avoid both by paying
every dense term as the change of a function of the state. Let $u(s) = 1 - f(S)$,
$k(s)=\lvert S\rvert$, and let
\[
  R(s) \;=\; \sum_{g_j \in S} \exp\!\bigl(-\hat\delta^-_j/\sigma_p\bigr), \qquad \sigma_p = 0.005,
\]
be a soft count of redundant guards: a guard that sees nothing on its own counts as $1$, one
that alone covers well over half a percent of the gallery counts as about $0$. Here
$\hat\delta^-_j$ is the raster estimate of the exclusive area, the same quantity the agent
sees in observation channel 4. Define
\[
  \Pot(s) \;=\; -w_{\mathrm{cov}}\,u(s) \;-\; w_{\mathrm{par}}\,k(s) \;-\; w_{\mathrm{prune}}\,R(s),
  \qquad
  F(s,s') = \Pot(s') - \Pot(s).
\]
$F$ is exactly the dense reward we pay. Its coverage and guard-count parts are
$+w_{\mathrm{cov}}\delta^+ - w_{\mathrm{par}}$ on $\ADD$ and $-w_{\mathrm{cov}}\delta^- +
w_{\mathrm{par}}$ on $\DEL$; its redundancy part $w_{\mathrm{prune}}(R(s)-R(s'))$ is paid on
every step. Along any episode,
\[
  \sum_{t=0}^{T-1} F(s_t,s_{t+1}) \;=\; \Pot(s_T) - \Pot(s_0)
  \;=\; w_{\mathrm{cov}}\,c_T \;-\; w_{\mathrm{par}}\,k_T \;-\; w_{\mathrm{prune}}\,R(s_T) \;+\; \text{const},
\]
where the constant depends only on the start state ($0$ from an empty gallery). The dense
reward therefore depends only on \emph{where the episode ends}, not on the path to it:
coverage reached, guards used, and how many of them are redundant. No sequence of actions
can collect more than that endpoint is worth, so no dense term can be farmed. The potentials
are differenced with $\gamma=1$ and the terminal state is not zeroed, so the endpoint value is
part of the return rather than cancelled; it is exactly the objective of \cref{sec:collapse}
(cover, with few guards, none of them redundant). The identity is exact for
the undiscounted return. PPO discounts with $\gamma=0.995$, and episodes are at most 40
steps long, so later terms are weighted by at least $0.995^{40}\approx0.82$: the shaping is
close to, though not exactly, the discounted form $\gamma\Pot(s')-\Pot(s)$ whose policy
invariance \citet{ng1999} prove.

\begin{designrule}[Anneal iff non-potential]\label{rule:anneal}
A reward term may keep a fixed weight for the whole of training only if it is a potential
difference or part of the objective. Any other term biases the optimum and must reach zero
before training ends.
\end{designrule}

\Cref{rule:anneal} decides the schedule in \cref{fig:curriculum}. The coverage,
guard-count and redundancy terms are potentials and keep their weights. The terminal terms
are part of the objective. The regularity penalties ($w_{\mathrm{proj}}$, $w_{\mathrm{void}}$,
$w_{\mathrm{stop}}$) only discourage actions that are never useful, so they cannot be farmed and
change no optimum. The stage-1 scaffolds (frontier, sweep, separation) are the only terms
that bias the solution, and they are annealed to zero.

\begin{figure}[!htbp]
\centering
\begin{tikzpicture}
\begin{axis}[
  width=0.96\linewidth, height=6.4cm,
  xmin=0, xmax=2.5, ymin=0, ymax=0.30,
  xlabel={training steps ($\times 10^6$)}, ylabel={reward weight},
  xtick={0,0.5,0.7,1,1.05,1.5,2}, xticklabels={0,0.5,0.7,1,,1.5,2},
  ytick={0,0.05,0.10,0.15,0.20},
  yticklabel style={/pgf/number format/fixed, /pgf/number format/precision=2},
  tick label style={font=\small}, label style={font=\small},
  legend style={font=\small, draw=none, fill=white, fill opacity=0.85, text opacity=1,
                at={(0.5,1.02)}, anchor=south, legend columns=5},
  every axis plot/.append style={thick}]
  % stage backgrounds
  \fill[agpblue!6]   (axis cs:0,0)   rectangle (axis cs:1,0.30);
  \fill[agpgreen!8]  (axis cs:1,0)   rectangle (axis cs:2,0.30);
  \fill[pattern=north east lines, pattern color=black!12] (axis cs:2,0) rectangle (axis cs:2.5,0.30);
  \addplot[agpblue]                 coordinates {(0,0.10) (0.7,0.10) (1,0) (2,0)};
  \addplot[agpblue, dashed]         coordinates {(0,0.06) (0.7,0.06) (1,0) (2,0)};
  \addplot[agpblue, dotted]         coordinates {(0,0.04) (0.7,0.04) (1,0) (2,0)};
  \addplot[agpaccent]               coordinates {(0,0) (1,0) (1.05,0.15) (2,0.15)};
  \addplot[agpaccent, dashed]       coordinates {(0,0) (1,0) (1.05,0.20) (2,0.20)};
  \legend{$w_{\mathrm{sep}}$, $w_{\mathrm{sweep}}$, $w_{\mathrm{front}}$, $w_{\mathrm{par}}$, $w_{\mathrm{prune}}$}
  \node[font=\small, align=center] at (axis cs:0.4,0.262) {\textbf{Stage 1: \stageI{}}\\only $\ADD$; guards free};
  \node[font=\small, align=center] at (axis cs:1.5,0.262) {\textbf{Stage 2: \stageII{}}\\$\ADD$, $\DEL$, $\STOP$; guards cost};
  \node[font=\small, align=center] at (axis cs:2.25,0.255) {\textbf{Stage 3}\\(optional)\\$+\MOVE$};
  \node[font=\footnotesize, align=center] at (axis cs:0.85,0.135) {scaffolds\\annealed};
  \node[font=\footnotesize, align=left, anchor=west] at (axis cs:1.08,0.045) {5\% blend; 40\% of episodes start\\from greedy cover $+\,4$ spares};
\end{axis}
\end{tikzpicture}
\caption{\textbf{The curriculum and its reward-weight schedule}, with the exact values from
\code{rewards.py}. In stage 1 the scaffolds (blue) hold their weight to $70\%$ of the stage
and then fall linearly to $0$, so the last $30\%$ trains on the unbiased reward
(\cref{rule:anneal}). Stage 2 blends from that column to the parsimony weights (orange) over
its first $5\%$, which spares the value function a step change. Not plotted because of scale:
the potential term $w_{\mathrm{cov}}=1$ and the terminal $w_{\mathrm{goal}}=4$,
$w_{\mathrm{fail}}=2$ are constant; $w_{\mathrm{pars}}$ ($0\to0.5$), $w_{\mathrm{stop}}$
($0\to2$) and $w_{\mathrm{void}}$ ($0\to0.1$) blend in with $w_{\mathrm{par}}$.}
\label{fig:curriculum}
\end{figure}

% ===========================================================================
\section{Architecture}\label{sec:arch}
% ===========================================================================

\subsection{What the policy has to do}

At every step the policy outputs a mode in $\{\ADD,\DEL,\STOP\}$ and a point $(x,y)$ in the
gallery's bounding box, normalised to $[0,1]^2$. The geometry of \cref{sec:prelim} turns this
into five concrete requirements, and each part of the network answers one of them:

\begin{description}[leftmargin=1.8em]
  \item[R1 Continuous placement.] Point guards are continuous (\cref{sec:complexity}), so
        the action must reach any point, not a fixed candidate set.
  \item[R2 Several targets at once.] A partly covered hard gallery usually has several
        separate uncovered pockets. The distribution over points must be able to put mass
        on all of them, rather than on their average.
  \item[R3 Global context.] With unlimited range, whether a point is useful depends on what
        it can see across the whole gallery, not on its neighbourhood.
  \item[R4 Fine localisation.] The missing $2$--$4\%$ of a nearly solved gallery is often a
        thin wedge behind a hole or a reflex corner (\cref{fig:visibility}), a few pixels
        wide. The network must not blur it away before deciding where to aim.
  \item[R5 Unordered modes with masking.] $\ADD$, $\DEL$ and $\STOP$ are unordered
        alternatives, and stage 1 must disable two of them cleanly.
\end{description}

\subsection{The network}

\Cref{fig:net} shows the network (\code{policy.SpatialPolicy}, about $294$k parameters). Its
organising idea is to \emph{read the placement off a map} rather than compute coordinates
from a summary vector: one branch keeps the image's layout down to a $16\times16$ grid of
cells, and a second branch summarises the whole gallery into one vector that is handed to
every cell. The mode and the value of the state, which are properties of the whole gallery,
are read from the summary.

\begin{figure}[!htbp]
\centering
\begin{tikzpicture}[node distance=0.55cm]
  \node[blk, text width=5.0cm] (in) {\textbf{Input} $64{\times}64{\times}7$\\5 image channels $+$ $x$, $y$ coordinates};
  \node[blk, text width=5.0cm, below=of in] (l1) {conv $5{\times}5$/2 $\;\to\; 32{\times}32{\times}32$};
  \node[blk, text width=5.0cm, below=of l1] (l2) {conv $3{\times}3$/2 $\;\to\; 16{\times}16{\times}32$};
  \node[blk, text width=5.0cm, below=of l2] (l3) {conv $3{\times}3$ $\;\to\;$ \textbf{cell features}\\$16{\times}16{\times}32$};
  \node[blk, text width=5.0cm, below=1.3cm of l3] (cat) {concat $\;\to\; 16{\times}16{\times}64$};
  \node[blk, text width=5.0cm, below=of cat] (c1) {\textbf{Cell head}\\$2\times$ conv $3{\times}3$ $\to 16{\times}16{\times}32$\\conv $1{\times}1$ $\to 16{\times}16{\times}5$};
  \node[blk, text width=5.0cm, below=of c1, fill=agpaccent!12, draw=agpaccent] (out) {per cell $k$: score $s_k$ and\\Beta parameters $a^x_k, b^x_k, a^y_k, b^y_k$};

  \node[blk, text width=4.6cm, right=1.2cm of l3] (g1) {conv $3{\times}3$/2 $\to 8{\times}8{\times}32$\\conv $3{\times}3$/2 $\to 4{\times}4{\times}32$};
  \node[blk, text width=4.6cm, below=of g1] (g2) {flatten $512$ $\to$ FC $256$\\concat 5 stats $\to$ FC $256$};
  \node[blk, text width=4.6cm, below=of g2, fill=agpblue!15] (z) {\textbf{global summary} $z\in\Real^{256}$};
  \node[blk, text width=4.6cm, below=of z] (bc) {\textbf{broadcast}: FC $\to 32$,\\copied to all $16{\times}16$ cells};
  \node[blk, text width=4.6cm, below=1.1cm of bc, fill=agpaccent!12, draw=agpaccent] (mode) {\textbf{Mode head}: FC $\to$ 3 logits\\$\ADD$ / $\DEL$ / $\STOP$};
  \node[blk, text width=4.6cm, below=of mode, fill=agpaccent!12, draw=agpaccent] (val) {\textbf{Value head}: FC 64, tanh, FC $\to V(s)$};

  \node[blk, text width=2.3cm, left=0.6cm of l3, fill=white] (mask) {interior mask\\max-pool ch.\ 0\\$\to 16{\times}16$};

  \draw[flow] (in) -- (l1); \draw[flow] (l1) -- (l2); \draw[flow] (l2) -- (l3);
  \draw[flow] (l3) -- (cat); \draw[flow] (cat) -- (c1); \draw[flow] (c1) -- (out);
  \draw[flow] (l3) -- (g1); \draw[flow] (g1) -- (g2); \draw[flow] (g2) -- (z);
  \draw[flow] (z) -- (bc); \draw[flow] (bc.west) -- ++(-0.6,0) |- (cat.east);
  \draw[flow] (z.east) -- ++(0.4,0) |- (mode.east);
  \draw[flow] (z.east) -- ++(0.4,0) |- (val.east);
  \draw[flow] (in.west) -| (mask.north);
  \draw[flow] (mask.south) |- (out.west);
  \node[font=\small\bfseries, text=agpblue, above=0.05cm of in] {local branch (keeps the layout)};
  \node[font=\small\bfseries, text=agpblue, above=0.05cm of g1] {global branch (whole gallery)};
\end{tikzpicture}
\caption{\textbf{The spatial policy, with tensor sizes.} Left: the local branch never
flattens the image, so placement is scored per cell on a $16\times16$ map. Right: the global
branch summarises the whole gallery and the scalar stats into $z$, which is broadcast back to
every cell (unlimited-range visibility depends on the whole gallery) and also feeds the mode
and value heads. The interior mask removes cells with no free space from the cell
distribution. ReLU after every convolution; ``/2'' is stride 2.}
\label{fig:net}
\end{figure}

\paragraph{A $16\times16$ map, never flattened (R4).}
The local branch downsamples the $64\times64$ image only to the cell grid ($G=16$, each cell
$4\times4$ pixels) and keeps it there. Every cell has its own feature vector, computed from
its own neighbourhood, so a thin uncovered wedge stays a feature of the cells it touches and
can raise their scores directly. Aiming at a corner is then no harder than aiming at open
floor: it is a different cell of the same map. Flattening the image into one vector before
choosing a position would force the network to convert ``there is a sliver at the top left''
into numeric coordinates, which is exactly the kind of fine spatial detail a coarse summary
loses.

\paragraph{Coordinate channels.}
Convolutions are translation-invariant, but visibility is not: a point next to a wall sees
differently from the same local pattern in the middle of the room. Two extra input channels
hold each pixel's normalised $x$ and $y$ \cite{liu2018coordconv}, so every cell knows where it
is in the gallery.

\paragraph{A global summary broadcast to every cell (R3).}
Local features alone cannot tell whether a point would see an uncovered region on the far
side of the gallery. The global branch continues to $4\times4$, flattens, and merges the
result with the five scalar stats into $z\in\Real^{256}$. A linear map of $z$ to $32$
numbers is copied to every cell and concatenated with that cell's local features, so each
cell is scored from both its neighbourhood and the state of the whole gallery: what is
covered, where the guards are, how much budget is left.

\paragraph{Masking cells outside the gallery.}
A cell is \emph{live} if any of its $4\times4$ raster samples lies in the free space (a max-pool
of the interior channel). Dead cells, outside $\Poly$ or entirely inside a hole, get logit
$-10^{8}$ and zero probability (\cref{fig:action}). This removes a large and useless part of
the search space (in the gallery of \cref{fig:action}, $68$ of the $256$ cells), and the
policy never wastes probability on positions that would only be projected back into the
gallery and penalised.

\paragraph{Heads and initialisation.}
The cell head outputs five maps: a score $s_k$ per cell and four Beta parameters, made
$\ge 1$ by $1+\operatorname{softplus}(\cdot)$ so that every offset density is unimodal and
bounded. The mode and value heads read $z$, because ``stop now'' and ``how good is this
state'' are global questions. The last layers of the cell and mode heads are initialised
orthogonally with gain $0.01$, so training starts from near-uniform cells, near-uniform modes
and gently centred offsets.

\subsection{The factored action}\label{sec:action}

Index cells by $k = rG + c$ with row $r$ along $y$ and column $c$ along $x$. The policy
samples
\[
  \mu \sim \mathrm{Cat}(\pi^{\mathrm{mode}}), \qquad
  k \sim \mathrm{Cat}(\pi^{\mathrm{cell}}), \qquad
  o_x \sim \mathrm{Beta}(a^x_k,b^x_k), \quad o_y \sim \mathrm{Beta}(a^y_k,b^y_k),
\]
\[
  (x,y) \;=\; \Bigl(\frac{c + o_x}{G},\; \frac{r + o_y}{G}\Bigr),
  \qquad
  \pi^{\mathrm{cell}}_k \propto \exp(s_k)\,\Ind{\text{cell $k$ is live}},
\]
and sends the vector $[\mu, x, y]$ to the environment (\cref{fig:action}).

\begin{figure}[!htbp]
\centering
\begin{tikzpicture}[x=0.5cm, y=0.5cm]
  % dead cells
  \foreach \c/\r in {0/0,1/0,2/0,3/0,4/0,5/0,6/0,7/0,14/0,15/0,0/1,1/1,2/1,3/1,4/1,15/1,0/2,1/2,2/2,3/2,0/3,1/3,2/3,3/3,0/4,1/4,2/4,0/5,1/5,2/5,8/5,9/5,0/6,1/6,8/6,0/7,0/8,7/8,13/8,6/9,7/9,13/9,7/10,0/11,0/12,1/12,0/13,1/13,2/13,3/13,0/14,1/14,2/14,3/14,4/14,14/14,15/14,0/15,1/15,2/15,3/15,4/15,5/15,6/15,12/15,13/15,14/15,15/15}
    \fill[masked] (\c,\r) rectangle ++(1,1);
  % gallery (real mid-tier test gallery #30, scaled to the 16x16 grid)
  \begin{scope}[x=1cm, y=1cm]
    \fill[seen, opacity=0.6] (7.750,4.006) -- (8.000,6.409) -- (5.952,7.700) -- (4.121,8.000) -- (2.784,7.173) -- (0.000,5.017) -- (1.540,2.775) -- (2.616,0.751) -- (4.596,0.310) -- (6.365,0.000) -- (7.999,1.325) -- cycle;
    \draw[outer] (7.750,4.006) -- (8.000,6.409) -- (5.952,7.700) -- (4.121,8.000) -- (2.784,7.173) -- (0.000,5.017) -- (1.540,2.775) -- (2.616,0.751) -- (4.596,0.310) -- (6.365,0.000) -- (7.999,1.325) -- cycle;
    \draw[hole] (3.240,4.141) -- (3.803,3.958) -- (4.297,4.212) -- (4.501,4.921) -- (4.174,5.535) -- (3.631,5.787) -- (3.152,5.499) -- (2.916,4.728) -- cycle;
    \draw[hole] (6.893,5.075) -- (6.448,5.056) -- (6.171,4.630) -- (6.285,4.041) -- (6.650,3.868) -- (7.059,4.104) -- (7.170,4.706) -- cycle;
    \draw[hole] (4.467,2.330) -- (4.995,2.557) -- (4.942,3.294) -- (4.551,3.654) -- (4.097,3.568) -- (3.711,3.078) -- (3.994,2.426) -- cycle;
  \end{scope}
  \draw[black!35, line width=0.3pt] (0,0) grid (16,16);
  \draw[black!70] (0,0) rectangle (16,16);
  % chosen cell (c=11, r=10) and the continuous point inside it
  \draw[agpaccent, line width=1.4pt] (11,10) rectangle ++(1,1);
  \node[guard, fill=agpaccent] (p) at (11.62,10.35) {};
  \node[font=\small, anchor=north] at (8,-0.2) {$x \in [0,1]$ \;($c = 0,\dots,15$)};
  \node[font=\small, rotate=90, anchor=south] at (-0.2,8) {$y \in [0,1]$ \;($r = 0,\dots,15$)};

  % zoomed cell
  \begin{scope}[shift={(19.5,6)}, x=3cm, y=3cm]
    \fill[seen, opacity=0.6] (0,0) rectangle (1,1);
    \draw[agpaccent, line width=1.4pt] (0,0) rectangle (1,1);
    \draw[dotted, black!60] (0.62,0) -- (0.62,0.35) -- (0,0.35);
    \draw[<->, black!75] (0.02,0.07) -- node[above, font=\small] {$o_x$} (0.62,0.07);
    \draw[<->, black!75] (0.75,0.02) -- node[right, font=\small] {$o_y$} (0.75,0.35);
    \node[guard, fill=agpaccent, label={[font=\small]above:$(x,y)$}] at (0.62,0.35) {};
    \node[font=\small, anchor=south] at (0.5,1.03) {cell $k = rG + c$, \;$(c,r) = (11,10)$};
    \node[font=\small, anchor=north, align=center] at (0.5,-0.08)
      {$(x,y) = \bigl(\tfrac{c+o_x}{G}, \tfrac{r+o_y}{G}\bigr)$\\$= (0.726,\, 0.647)$};
  \end{scope}
  \draw[black!50, densely dashed] (12,11) -- (19.5,9);
  \draw[black!50, densely dashed] (12,10) -- (19.5,6);
\end{tikzpicture}
\caption{\textbf{A region is chosen discretely, the point inside it continuously.} Left: the
$16\times16$ cell grid over the bounding box of a real gallery (mid test gallery \#30). Grey
cells contain no free space and are masked. The policy samples a live cell (orange, $(c,r)=(11,10)$),
then a continuous offset from that cell's two Beta distributions,
$o_x\sim\mathrm{Beta}(a^x_k,b^x_k)$ and $o_y\sim\mathrm{Beta}(a^y_k,b^y_k)$. Right: the offset fixes
the exact point, here $o_x=0.62$, $o_y=0.35$. Any point in the gallery can be produced, so the
grid organises the choice without restricting where a guard can go.}
\label{fig:action}
\end{figure}

\paragraph{Why a distribution over regions (R2).}
When several parts of the gallery are still uncovered, the policy should be able to aim at
any of them. A categorical over cells assigns independent probability to every region, so it
can put mass on two opposite corners at once. A single Gaussian over $(x,y)$, or any method
that outputs one averaged point, cannot: its mean falls between the targets, usually in space
that is already covered (\cref{fig:averaging}). The same holds for the deterministic action
used in evaluation: the most probable cell is an actual choice of one region, while the mean
of a unimodal distribution is a compromise between regions.

\begin{figure}[!htbp]
\centering
\begin{tikzpicture}[x=0.32cm, y=0.32cm]
  \begin{scope}
    \fill[seen] (0,0) rectangle (16,16);
    \fill[unseen] (0.5,12) rectangle (4,15.5);
    \fill[unseen] (12,0.5) rectangle (15.5,4);
    \draw[outer] (0,0) rectangle (16,16);
    \foreach \s in {1.0,2.0,3.0} \draw[agpblue, line width=0.8pt] (8,8) ellipse ({\s*2.4} and {\s*2.4});
    \node[star, star points=5, fill=agpaccent, draw=black, inner sep=1.6pt] at (8,8) {};
    \node[font=\small, align=center, fill=white, fill opacity=0.8, text opacity=1] at (8,5.0) {mean lands in\\covered space};
    \node[font=\small, anchor=south] at (8,16.3) {\textbf{one Gaussian over $(x,y)$}};
    \node[font=\footnotesize] at (2.25,13.75) {pocket};
    \node[font=\footnotesize] at (13.75,2.25) {pocket};
  \end{scope}
  \begin{scope}[shift={(21,0)}]
    \fill[seen] (0,0) rectangle (16,16);
    \fill[unseen] (0.5,12) rectangle (4,15.5);
    \fill[unseen] (12,0.5) rectangle (15.5,4);
    \foreach \c/\r/\o in {0/12/45,1/12/70,2/12/70,3/12/45,0/13/70,1/13/100,2/13/100,3/13/70,0/14/70,1/14/100,2/14/100,3/14/70,0/15/45,1/15/70,2/15/70,3/15/45,
                          12/0/45,13/0/70,14/0/70,15/0/45,12/1/70,13/1/100,14/1/100,15/1/70,12/2/70,13/2/100,14/2/100,15/2/70,12/3/45,13/3/70,14/3/70,15/3/45}
      \fill[agpaccent!\o] (\c,\r) rectangle ++(1,1);
    \draw[black!25, line width=0.2pt] (0,0) grid (16,16);
    \draw[outer] (0,0) rectangle (16,16);
    \node[font=\small, align=center, fill=white, fill opacity=0.8, text opacity=1] at (8,8) {probability on\\both pockets};
    \node[font=\small, anchor=south] at (8,16.3) {\textbf{categorical over cells}};
  \end{scope}
\end{tikzpicture}
\caption{\textbf{Why the policy chooses a region instead of one averaged point.} Two
uncovered pockets (orange) sit in opposite corners of an otherwise covered gallery
(schematic). \emph{Left:} a single Gaussian that tries to serve both pockets centres between
them, so its mean, the deterministic action, lands in already-covered space, and its samples
mostly do too. \emph{Right:} the categorical over cells puts high probability (darker cells) on
both pockets independently, so every sample aims at real work, and the most probable cell is
one of the pockets.}
\label{fig:averaging}
\end{figure}

\paragraph{Why a continuous offset inside the cell (R1).}
The Beta offset makes the action continuous while keeping the density exact.

\begin{proposition}[The action is continuous]\label{prop:continuous}
For $\mu\neq\STOP$, the induced distribution of $(x,y)$ has density
\[
  p(x,y) \;=\; G^2\;\pi^{\mathrm{cell}}_{k(x,y)}\;
  \mathrm{Beta}\bigl(o_x(x)\bigr)\,\mathrm{Beta}\bigl(o_y(y)\bigr)
\]
on $[0,1]^2$, where $k(x,y) = \lfloor Gy\rfloor G + \lfloor Gx\rfloor$,
$o_x(x) = Gx - \lfloor Gx\rfloor$ and $o_y(y) = Gy - \lfloor Gy\rfloor$. It is positive on every
live cell.
\end{proposition}

\begin{proof}
Cells partition $[0,1)^2$, so each point belongs to exactly one cell and determines its
offsets uniquely. Within a cell the map from offsets to $(x,y)$ scales each coordinate by
$1/G$, so the Jacobian contributes $G^2$. Beta densities with parameters $\ge1$ are positive on
$(0,1)$.
\end{proof}

The grid organises \emph{how} the policy decides; it does not restrict \emph{where} a guard
can go. This is the difference from candidate-set solvers \cite{couto2011,tozoni2016} and
from a plain discrete ``spatial action map'' \cite{wu2020spatial}, which would snap every
guard to a cell centre and give up advantage (1) of \cref{sec:complexity}. The Beta
distribution is the natural choice for a bounded offset: it lives on $(0,1)$ exactly, so no
probability is clipped at the cell edges, which avoids the boundary bias of clipped Gaussians
\cite{chou2017beta}.

\paragraph{Log-probability, entropy, deterministic action.}
Given a stored action $[\mu,x,y]$, the code recovers $k$, $o_x$, $o_y$ and computes
\[
  \log\pi(a\mid s) \;=\; \log\pi^{\mathrm{mode}}_\mu
  \;+\; \Ind{\mu\neq\STOP}\Bigl(\log\pi^{\mathrm{cell}}_k + \log\mathrm{Beta}(o_x) + \log\mathrm{Beta}(o_y)\Bigr).
\]
The constant $\log G^2$ is dropped because it cancels in PPO's probability ratio. The point
term is excluded for $\STOP$ because a $\STOP$ ignores the point, so including it would only add
gradient noise to a decision that does not depend on it. Offsets are clamped to
$[10^{-5}, 1-10^{-5}]$, which keeps sampled points strictly inside their cell and the Beta
log-densities finite. The entropy bonus is that of the joint choice,
\[
  \mathcal{H} \;=\; \mathcal{H}(\pi^{\mathrm{mode}}) + \mathcal{H}(\pi^{\mathrm{cell}})
  + \textstyle\sum_k \pi^{\mathrm{cell}}_k\bigl(\mathcal{H}(\mathrm{Beta}^x_k) + \mathcal{H}(\mathrm{Beta}^y_k)\bigr).
\]
The deterministic action used for evaluation is the most probable mode, the most probable
cell, and the Beta means $a/(a+b)$ inside it.

\subsection{The mode head across the stage switch (R5)}

\paragraph{A categorical, not a number.}
The three modes are separate logits. Moving probability from $\ADD$ to $\STOP$ does not pass
through $\DEL$, as it would if the mode were one continuous output cut into bands: learning
``stop once covered'' must not require first becoming worse at something else.

\paragraph{Stage 1: masking at the logits.}
The $\DEL$ and $\STOP$ logits are replaced by $-10^{8}$, so their probability is exactly zero
and they contribute nothing to the log-probability or the entropy \cite{huang2022masking}.
The environment also decodes every stage-1 action as $\ADD$, so the MDP is correct for any
policy.

\paragraph{Stage 2: priming.}
Masked logits receive no gradient, so at the switch the mode head is still at its
initialisation and would choose among the three modes arbitrarily. Instead, we zero the mode
head's weights and set its biases so that $\pi^{\mathrm{mode}} = (0.8, 0.1, 0.1)$ in every state
(\code{--add-prob 0.8}). Stage 2 therefore starts by behaving as stage 1 ended, mostly
adding, while sampling $\DEL$ and $\STOP$ often enough to learn their value. The placement
heads are left as stage 1 trained them.

\subsection{PPO configuration}

\begin{table}[!htbp]
\centering\small
\caption{PPO settings (\code{train.py}). One model runs both stages; the step counter
continues across the switch.}
\label{tab:ppo}
\begin{tabular}{ll@{\hspace{2.5em}}ll}
\toprule
learning rate & $3\times10^{-4}$ (stage 2: $10^{-4}$) & discount $\gamma$ & $0.995$ \\
environments & 16 (subprocesses) & GAE $\lambda$ & $0.95$ \\
rollout length & 256 per env (4096 per update) & clip range & $0.2$ \\
minibatch & 512 & epochs per update & 10 (stage 2: 5) \\
entropy coefficient & $0.005$ (both stages) & value coefficient & $0.5$ \\
gradient clip & $0.5$ & target KL & off \\
steps & $10^6$ per stage & device & Apple MPS \\
\bottomrule
\end{tabular}
\end{table}

\Cref{tab:ppo} lists the settings. $\gamma = 0.995$ is close to $1$ because episodes are short
(at most 40 steps) and the terminal bonus must reach the first placement. We do not stop
updates early on a KL threshold: a cap of $0.03$ ended updates after so few gradient steps
that the placement head barely trained.

\paragraph{Smaller steps in stage 2.}
Stage 2 changes the reward under a policy that already covers every tier. Its job is to
add $\DEL$ and $\STOP$ on top of that skill, not to relearn placement, so it uses a third of
stage 1's learning rate and half the epochs per update. At stage 1's step size, single
updates reached a KL divergence above $1$ late in stage 2, and the placement skill on hard
galleries degraded within $150$k steps (solved fell from about $95\%$ to $50\%$ in training).

\paragraph{No extra entropy bonus.}
The bonus acts on the joint entropy of \cref{sec:action}, and that is dominated by the
cell: up to $\ln 256 \approx 5.5$ nats against $\ln 3 \approx 1.1$ for the mode. Raising it to
encourage the new modes would mainly flatten the cell distribution, i.e.\ make placement
less precise, and hard galleries need precise corner guards. Exploration of $\DEL$ and
$\STOP$ comes from the priming instead, which samples them $20\%$ of the time at the switch.

\paragraph{Model selection.}
Even with smaller steps, the last stage-2 policy need not be the best one. Every $100$k
steps we save a checkpoint and score it on a separate validation set (60 galleries, 20 per
tier, generated from their own seed; \cref{tab:tiers}), running each gallery once with
deterministic actions. A gallery scores
\[
  q = \begin{cases} k_{\mathrm{ref}} / k & \text{if solved,}\\ 0 & \text{otherwise,}\end{cases}
\]
so $q = 1$ matches greedy and $q > 1$ beats it; the checkpoint's score is the mean of $q$, and
the highest-scoring checkpoint is the result. The starting policy is scored too, so the kept
model is never worse on validation than where stage 2 began. One number is needed because
ranking by solved rate first fails: a policy that never stops fills the 16-guard budget,
covers everything, and would outrank one that solves slightly fewer galleries with a third
of the guards. The test set is used once, for the final numbers in \cref{tab:results}.

\paragraph{Hardware.}
Training runs on Apple's MPS GPU: a forward and
backward pass over 512 samples takes $12$\,ms there versus $600$\,ms on the CPU. Beta sampling,
the one operation MPS lacks, is done on the CPU for the single offset drawn per environment
per step. $G=16$ is a compromise: $256$ regions separate distinct pockets while keeping the
categorical small enough to explore, and $G=32$ would double the network's cost without
changing what positions are reachable.

\subsection{Results}\label{sec:results}

\takeaway{%
\textbf{Results pending.} Stage 2 with the spatial policy (\code{runs/agp\_v3}, continuing
from the stage-1 checkpoint of \code{runs/agp\_v2}) was in progress when this report was
written. \Cref{tab:results} is a placeholder to be filled from the test-set evaluation of
the best-on-validation checkpoint (deterministic actions, 90 test galleries).}

\begin{table}[!htbp]
\centering\small
\caption{Test-set evaluation of the spatial policy (best stage-2 checkpoint on validation).
\pending{: fill from \code{runs/train\_v3.log}}}
\label{tab:results}
\begin{tabular}{lcccc}
\toprule
tier & solved & mean coverage & mean $k$ (solved) & $k/k_{\mathrm{ref}}$ (solved) \\
\midrule
easy & \pending{} & \pending{} & \pending{} & \pending{} \\
mid  & \pending{} & \pending{} & \pending{} & \pending{} \\
hard & \pending{} & \pending{} & \pending{} & \pending{} \\
\bottomrule
\end{tabular}
\end{table}

% ===========================================================================
\section{MDP formulation and reward}\label{sec:mdp}
% ===========================================================================

\Cref{fig:mdp} summarises one step of the decision process. The rest of this section states
each part precisely.

\begin{figure}[!htbp]
\centering
\begin{tikzpicture}
  \node[blk, text width=5.6cm] (state) at (0,0) {\textbf{State} $s_t$\\gallery, guards $S_t$, step $t$,\\best coverage so far $c^{\mathrm{best}}_t$};
  \node[blk, text width=5.6cm] (obs) at (7.4,0) {\textbf{Observation}\\image $64{\times}64{\times}5$: interior, covered,\\uncovered, guards, redundancy\\5 stats: coverage, budget, time,\\cheapest deletion, ratchet room};
  \node[blk, text width=5.6cm] (pol) at (7.4,-3.0) {\textbf{Policy} (\cref{fig:net})\\action $a_t = [\mu, x, y]$};
  \node[blk, text width=5.6cm, fill=white] (env) at (0,-3.0) {\textbf{Environment step}\\$\ADD$: project into $\Free$, place a guard\\$\DEL$: remove the nearest guard,\\unless the ratchet refuses\\$\STOP$: end the episode};
  \node[blk, text width=5.6cm, fill=agpaccent!12, draw=agpaccent] (rew) at (0,-5.9) {\textbf{Reward} $r_t$ (exact geometry)\\coverage, guard and redundancy potentials,\\scaffolds, penalties, terminal bonus};
  \draw[flow] (state) -- node[above, font=\small]{render} (obs);
  \draw[flow] (obs) -- (pol);
  \draw[flow] (pol) -- node[above, font=\small]{$a_t$} (env);
  \draw[flow] (env) -- (rew);
  \draw[flow] (env.west) -- ++(-0.5,0) |- node[pos=0.25, left, font=\small]{$s_{t+1}$} (state.west);
  \draw[flow] (rew.east) -| node[pos=0.25, below, font=\small]{to PPO} (pol.south);
\end{tikzpicture}
\caption{\textbf{One step of the MDP.} The state is rendered into a 5-channel image and five
scalars; the policy outputs a mode and a point; the environment applies it with exact
geometry and returns the next state and a reward. Transitions are deterministic; the only
randomness is the gallery drawn at reset and the policy's own sampling.}
\label{fig:mdp}
\end{figure}

\subsection{Instances}

Each episode draws a gallery uniformly from the training set and applies a random symmetry
of the square (one of 8 rotations and reflections about the gallery centre), which also
transforms the stored greedy guards. Augmentation is switched off for evaluation. The
generator (\code{polygons.py}) produces three tiers (\cref{tab:tiers}).

\begin{table}[!htbp]
\centering\small
\caption{Difficulty tiers. 200 training, 20 validation and 30 test galleries per tier, each
split from its own seed.}
\label{tab:tiers}
\begin{tabular}{lccc}
\toprule
tier & outer vertices & holes & vertices per hole \\
\midrule
easy & 8--12  & 1--3 & 4--6 \\
mid  & 10--16 & 2--5 & 4--8 \\
hard & 16--26 & 5--9 & 5--10 \\
\bottomrule
\end{tabular}
\end{table}

\subsection{State, actions, termination}

The state is $s_t = (\Gallery, S_t, t, c^{\mathrm{best}}_t)$: the gallery, the multiset of
placed guards, the step count, and the best coverage reached so far in the episode. The action
$[\mu, x, y]$ is decoded to the point $q = (x_{\min} + x\,w,\; y_{\min} + y\,h)$.

\begin{description}[leftmargin=1.6em]
  \item[$\ADD$.] If $q\notin\operatorname{int}\Free$ (outside $\Poly$ or inside a hole), it is
        projected to the nearest boundary point and then moved towards a fixed interior
        reference point by the first $t\in\{0.01,0.05,0.10,0.25,0.50\}$ that lands strictly
        inside. A guard is \emph{always} placed: rejecting the action instead would leave the
        state unchanged, and a deterministic policy that emits the same invalid point would
        repeat it forever. The projection is penalised (\cref{sec:reward}) so the policy
        still learns the shape.
  \item[$\DEL$.] Removes the guard nearest to $q$. Addressing a guard by a point rather than
        an index keeps the action fixed-size and independent of guard order, and reuses the
        same point head as $\ADD$. Deletions are subject to a \emph{coverage ratchet}: a
        deletion is refused (state unchanged, penalised) if it would bring coverage below
        $c^{\mathrm{best}}_t - 0.01$. Undoing real progress is then impossible rather than
        merely costly.
  \item[$\STOP$.] Ends the episode with the current guard set.
\end{description}

\paragraph{Termination and truncation.}
The episode \emph{terminates} on $\STOP$; when an $\ADD$ reaches the budget of $16$ guards;
and, in stage 1 only, as soon as $f\ge c^\dagger$ (in stage 2 it continues so that pruning
can happen). It is \emph{truncated} after $40$ steps. Truncation is a safety net, not part of
the problem: no terminal reward is paid, and PPO bootstraps from the value function instead,
which would double-count if we also paid out.

\paragraph{Why a budget of 16.}
Greedy needs at most $11$ guards on any training gallery, so $16$ leaves room for stage 1 to
over-fill and stage 2 to prune. It is tight enough that guard placement still matters, and
the observation reports the fraction of budget used. The reward sizes in \cref{sec:calib}
are checked against this value and must be rechecked if it changes.

\subsection{Observation}\label{sec:obs}

\begin{table}[!htbp]
\centering\small
\caption{Observation. Image channels are $0/255$ masks except channel 4.}
\label{tab:obs}
\begin{tabularx}{\linewidth}{lX}
\toprule
\multicolumn{2}{l}{\emph{Image channels} ($64\times64$ over the bounding box)} \\
0 interior & pixels inside $\Free$ \\
1 covered & pixels seen by at least one guard \\
2 uncovered & interior pixels seen by no guard: the work left \\
3 guards & each guard as a $3\times3$ blob, so it survives strided convolutions \\
4 redundancy & each guard's blob shaded by $\min(1, \hat\delta^-_j/\sigma_p)$, where
$\hat\delta^-_j$ is the raster estimate of its exclusive area and $\sigma_p = 0.005$; dark
means nearly free to delete \\
\midrule
\multicolumn{2}{l}{\emph{Scalars} (all clipped to $[0,1]$)} \\
coverage & exact $f(S_t)$ \\
budget used & $k_t/16$ \\
time used & $t/40$ \\
cheapest deletion & $\min(1, \min_j \hat\delta^-_j/\sigma_p)$; $0$ means a free deletion exists now \\
ratchet room & $f(S_t) - (c^{\mathrm{best}}_t - 0.01)$: coverage a deletion may still cost \\
\bottomrule
\end{tabularx}
\end{table}

\Cref{tab:obs} lists the observation. Channel 2 tells the policy where the work is; channel 4
is what makes $\DEL$ learnable at all, since without it a load-bearing guard and a redundant
one look identical. The raster estimate $\hat\delta^-_j$ counts pixels seen by guard $j$ and
by no other guard, one vectorised operation per guard. The exact $\delta^-$, a polygon union
over the other guards, is computed only when a deletion actually happens, for the reward and
the ratchet.

\subsection{Reward}\label{sec:reward}

Let $g$ be the placed (projected) point, $\iota = \Ind{q \notin \operatorname{int}\Free}$,
$D$ the bounding-box diagonal and $j^\star$ the guard nearest $q$. The per-step reward is
\[
r_t \;=\; r^{\mu}_t \;+\; w_{\mathrm{prune}}\bigl(R(s_t) - R(s_{t+1})\bigr),
\]
\[
r^{\mu}_t \;=\;
\begin{cases}
  w_{\mathrm{cov}}\,\delta^+ - w_{\mathrm{par}} - w_{\mathrm{proj}}\,\iota + r^{\mathrm{scaf}}_t, & \mu=\ADD, \\[4pt]
  -w_{\mathrm{cov}}\,\delta^-_{j^\star} + w_{\mathrm{par}}, & \mu=\DEL\text{, accepted}, \\[4pt]
  -w_{\mathrm{void}}, & \mu=\DEL\text{, refused or no guards}, \\[4pt]
  -w_{\mathrm{stop}}\Ind{f(S_t) < c^\dagger}, & \mu=\STOP,
\end{cases}
\]
where $R$ is the soft redundant-guard count of \cref{sec:shaping}, recomputed after every
step whatever the mode, plus, on termination only,
\[
  r^{\mathrm{term}} \;=\;
  \begin{cases}
    w_{\mathrm{goal}} + w_{\mathrm{pars}}\max(k_{\mathrm{ref}} - k_T,\, 0), & f(S_T)\ge c^\dagger,\\
    -w_{\mathrm{fail}}, & \text{otherwise.}
  \end{cases}
\]
where the scaffold reward of a placement is
\[
  r^{\mathrm{scaf}}_t \;=\; w_{\mathrm{front}}\Ind{g\notin\Cov(S_t)} + w_{\mathrm{sweep}}\,\kappa_t(g)
  - w_{\mathrm{sep}}\Ind{\textstyle\min_j\lVert g-g_j\rVert < 0.15\,D}.
\]
The sweep kernel favours the wall early and the interior later:
$\kappa_t(g) = \exp\bigl(-(\dbd(g)/D - m_t)^2/(2\cdot0.05^2)\bigr)$ with
$m_t = 0.02 + 0.23\,k_t/16$. \Cref{tab:weights} gives every weight;
\cref{fig:curriculum} shows how they move over training.

\begin{table}[!htbp]
\centering\small
\caption{Reward terms and their weights (\code{rewards.py}). ``Potential'' marks the terms
that are differences of the potential in \cref{sec:shaping}.}
\label{tab:weights}
\begin{tabularx}{\linewidth}{llXcc}
\toprule
group & weight & role & \stageI{} & \stageII{} \\
\midrule
coverage   & $w_{\mathrm{cov}}$   & area gained / lost (potential) & 1.0 & 1.0 \\
           & $w_{\mathrm{goal}}$  & terminal bonus for $f\ge c^\dagger$ & 4.0 & 4.0 \\
           & $w_{\mathrm{fail}}$  & terminal penalty otherwise & 2.0 & 2.0 \\
parsimony  & $w_{\mathrm{par}}$   & cost per guard, refunded on deletion (potential) & 0 & 0.15 \\
           & $w_{\mathrm{prune}}$ & drop in the soft count of redundant guards $R$ (potential) & 0 & 0.20 \\
           & $w_{\mathrm{pars}}$  & terminal bonus per guard saved vs.\ $k_{\mathrm{ref}}$ & 0 & 0.50 \\
scaffolds  & $w_{\mathrm{front}}$ & guard placed in uncovered space & 0.04 & 0 \\
           & $w_{\mathrm{sweep}}$ & boundary-first, then inward & 0.06 & 0 \\
           & $w_{\mathrm{sep}}$   & penalty for a guard within $0.15D$ of another & 0.10 & 0 \\
regularity & $w_{\mathrm{proj}}$  & aimed outside $\Free$ & 0.25 & 0.25 \\
           & $w_{\mathrm{void}}$  & impossible or refused deletion & 0 & 0.10 \\
           & $w_{\mathrm{stop}}$  & $\STOP$ before $c^\dagger$ & 0 & 2.0 \\
\bottomrule
\end{tabularx}
\end{table}

\paragraph{Why each group exists.}
The coverage, guard-count and redundancy potentials carry the objective
(\cref{sec:shaping}); the redundancy potential is what pays for deleting redundant guards. The terminal
terms make crossing $c^\dagger$ the thing that matters, and $w_{\mathrm{pars}}$ rewards beating
greedy, never punishing a match. The scaffolds make early coverage learning fast: aim at
uncovered space, trace the boundary first, do not stack guards; being non-potential, they are
removed before stage 1 ends. The regularity terms keep actions meaningful: aim inside the
gallery, do not request impossible deletions, do not stop short.

\subsection{Schedule}\label{sec:schedule}

\begin{itemize}
  \item \textbf{Stage 1.} The scaffolds hold full weight for the first $70\%$ of the stage
        and then fall linearly to $0$ at its end (\cref{fig:curriculum}).
  \item \textbf{Stage 2.} The weights blend linearly from the stage-1 column (scaffolds at
        zero) to the stage-2 column over the first $5\%$ of the stage, so the value function
        is not hit by a step change in its targets. $40\%$ of stage-2 episodes start
        \emph{presaturated}: the stored greedy cover plus $4$ guards at random interior
        points. Those extras are redundant by construction, so free deletions exist from the
        first step and $\DEL$ gets a dense learning signal.
  \item A callback pushes the current weights to every environment every 4096 steps; the
        environment has no notion of stage beyond its weights and the \code{allow\_delete}
        flag.
\end{itemize}

\subsection{Weight calibration}\label{sec:calib}

The weights are sized so that no shortcut beats solving the gallery.
\begin{itemize}
  \item \textbf{No scaffold farming (stage 1).} $w_{\mathrm{front}} + w_{\mathrm{sweep}} \le
        0.10$ per placement, so a full budget of 16 earns at most $1.6$ in scaffolds, well
        below the $4$ for reaching $c^\dagger$ plus the $2$ for avoiding failure.
  \item \textbf{Parsimony never makes failure attractive (stage 2).} A full budget costs
        $16\times0.15 = 2.4$, less than the $w_{\mathrm{goal}}+w_{\mathrm{fail}} = 6$ swing
        between solving and failing.
  \item \textbf{Pruning is worth doing.} Deleting a fully redundant guard lowers $R$ by at
        least $1$, so it returns at least $w_{\mathrm{par}} + w_{\mathrm{prune}} = 0.35$, large
        relative to per-step noise. It returns more when the deletion also makes a partner
        load-bearing: if two guards cover the same region, both count as redundant, and
        removing one drops $R$ by $2$, for $0.15 + 0.20\times2 = 0.55$. A guard's contribution
        to $R$ decays as $\exp(-\hat\delta^-/0.005)$, by a factor $e$ for every half percent of
        the gallery it alone covers, so deleting a load-bearing guard earns almost nothing from
        $R$ and pays the coverage loss in full.
  \item \textbf{Stopping early is costly.} $\STOP$ below $c^\dagger$ costs $w_{\mathrm{stop}} =
        2$ plus the $w_{\mathrm{fail}} = 2$ terminal penalty.
\end{itemize}

\paragraph{Why the pruning reward is a potential.}
The obvious alternative, a bonus $w_{\mathrm{prune}}\exp(-\delta^-/\sigma_p)$ paid on each
deletion, can be farmed. Adding a guard in already covered space costs $w_{\mathrm{par}} =
0.15$, and deleting it straight away pays $w_{\mathrm{par}} + w_{\mathrm{prune}} = 0.35$: a net
$+0.20$ per two steps, which over a 40-step episode approaches the goal bonus. Paying the
change of $R$ instead closes the loop exactly: adding a redundant guard raises $R$ by $1$
and costs $w_{\mathrm{prune}}$, deleting it lowers $R$ by $1$ and pays $w_{\mathrm{prune}}$
back, and the cycle nets zero. Over any episode the pruning reward totals
$w_{\mathrm{prune}}(R(s_0)-R(s_T))$, which is bounded and is earned only by ending with fewer
redundant guards than at the start.

% ===========================================================================
\section{Optional stage 3: relocation}\label{sec:stage3}
% ===========================================================================

Stage 2 can only remove guards that are already redundant. It cannot move a guard so that a
\emph{neighbour} becomes redundant, which is where the remaining gap to the optimum lies, and
which greedy, never revisiting a choice, cannot do either. A third stage would add a $\MOVE$
mode: select the guard nearest to the chosen point and displace it by a bounded continuous
vector. Delete-then-add cannot substitute for it, because the ratchet refuses to delete any
guard that alone covers more than $1\%$ of the gallery.

What it would need:
\begin{itemize}
  \item \textbf{A displacement head.} The spatial policy accommodates this naturally: a second
        pair of Beta maps per cell gives a bounded displacement for the guard selected in that
        cell, and the log-probability factorises as in \cref{sec:action}.
  \item \textbf{Reward.} The potentials already value a move correctly: coverage pays
        $w_{\mathrm{cov}}(f(S') - f(S))$, the redundancy potential pays
        $w_{\mathrm{prune}}(R(s)-R(s'))$ when the move makes a neighbour redundant, and the
        guard count does not change. No new reward term is needed.
  \item \textbf{Ratchet interaction.} A move that lowers coverage below
        $c^{\mathrm{best}}-0.01$ should be refused like a deletion, or the slack widened for
        moves, since a good relocation can pass through a temporary loss.
  \item \textbf{Constraint-based regularity.} The curriculum design proposes moving the
        regularity terms into a per-step constraint solved with N-P3O or CPO
        \cite{lee2023constrained,achiam2017cpo}, with threshold $h = 1.5\times$ the greedy
        solver's mean per-step regularity cost. A fixed penalty is negligible next to the
        first guard's gain but dominant next to the last sliver's, which suppresses exactly
        the delicate end-game moves stage 3 is for.
  \item \textbf{Hard-instance sampling.} Reweight the tiers towards hard galleries and sample
        instances in proportion to their recent failure rate.
  \item \textbf{A stronger reference.} Claiming a gap to the optimum rather than to greedy
        needs an exact or near-exact solver \cite{couto2011,tozoni2016} on a sample of
        galleries.
\end{itemize}

The prerequisite is that stage 2 prunes reliably. Relocation exists to remove guards that
pruning cannot; if pruning does not yet work, adding $\MOVE$ doubles the continuous action
dimension without addressing the actual bottleneck.

\bibliographystyle{plainnat}
\bibliography{refs}

\end{document}
